% ============================================================
% section1_building.tex — Section I: Building the Team
% ============================================================

\chapter{Building the Team}

\sectionintro{``Before you can lead a frontier team, you have to build one.
That means cultivating capability before optimizing for output---and removing
every layer of friction that slows people down.''} 

The principles in this section address the foundational act of team construction.
Not hiring alone, but the deliberate shaping of capability, communication architecture,
and organizational clarity that makes frontier work possible in the first place.

% ----------------------------------------------------------
\section{Cultivate Broad Capability First}

\principle{Build team depth before optimizing for output.}

Before a frontier team can move fast, it needs range. The instinct in high-pressure
R\&D environments is to optimize immediately---assign specialists to their narrow lane
and drive throughput. This instinct is wrong at the frontier. Narrow specialization
before broad capability produces teams that are brittle: fast on known problems,
helpless on unknown ones.

Cultivating broad capability means deliberately exposing team members to adjacent
domains, encouraging cross-functional literacy, and investing in the generalist
foundation before layering in specialization. This feels slow. It pays back
exponentially when the problem space shifts---and at the frontier, it always does.

\subsection{Story}
\tbd[Story to be attached. Candidate: early team-building experience where
breadth-first investment paid off during an unexpected pivot.]

% ----------------------------------------------------------
\section{Remove the Layers}

\principle{Eliminate every bureaucratic layer that slows frontier work.}

Bureaucratic layers are not neutral. They have a cost measured in latency, in
decision cycles, in the quiet erosion of the team's sense of ownership. At the
frontier, where speed of learning is the primary competitive variable, that cost
is fatal.

Removing the layers is not about creating chaos. It is about being ruthlessly
honest about which layers add value and which ones merely add process. Most
organizations accumulate the latter without noticing. The leader's job is to
notice---and to cut.

This principle is the operational expression of ``Vision First.'' A team that
understands the full trade space can self-direct. Layers exist, in part, because
leaders don't trust that understanding. Build the understanding; remove the layers.

\subsection{Story}
\tbd[Story to be attached. Consider: the ``too many architects'' story as a
counter-example---perfection-over-progress via excessive architecture review.
Confirm placement with Eric before attaching.]

% ----------------------------------------------------------
\section{Altitude of Decomposition}

\principle{Use decomposition level as a framework for seniority and growth.}

How finely someone can decompose a problem is a reliable signal of their technical
and leadership maturity. Junior contributors see tasks. Mid-level contributors see
workstreams. Senior contributors see systems. Principal contributors see the
constraints that define which systems are even possible.

The altitude at which someone naturally operates---how far they zoom out before
the picture becomes coherent for them---tells you more about their readiness for
expanded responsibility than any performance review. Use it that way. Assign
problems at one altitude above where someone currently operates comfortably.
That is the growth edge.

\subsection{Story}
\tbd[Story to be attached.]

% ----------------------------------------------------------
\section{Stakeholder Registry}

\principle{Design your stakeholder communication as architecture, not as a contact list.}

A stakeholder registry is not a list of names and email addresses. It is a
communication architecture: who needs what information, at what cadence, at what
level of detail, and in what format. Built deliberately, it is one of the most
powerful tools a leader has for reducing organizational friction.

Most leaders treat stakeholder communication reactively---answering questions as
they come, updating when pressed. This creates noise for some stakeholders and
information vacuums for others. The registry flips this. You define the
communication contract in advance, and you deliver against it.

\subsection{Story}
\tbd[Story to be attached.]

% ----------------------------------------------------------
\section{Project vs.\ Operations Clarity}

\principle{Know which mode you are in. Project and operations require different leadership.}

Project mode and operations mode are not points on a spectrum. They are fundamentally
different organizational states that require different leadership instincts, different
metrics, and different team structures. The failure to distinguish them is one of the
most common causes of R\&D dysfunction.

In project mode, the goal is learning and delivery against an uncertain target.
Variance is expected; speed of adaptation is the metric. In operations mode, the
goal is reliable execution against a known target. Variance is a defect; consistency
is the metric. Leaders who apply operations instincts to project-mode teams slow
them down. Leaders who apply project instincts to operations-mode teams create chaos.

Know which mode you are in. Lead accordingly.

\subsection{Story}
\tbd[Story to be attached. Reference: Shenhar \& Dvir Diamond Model for
matching management approach to project type.]

% ----------------------------------------------------------
\section{The Junkyard Principle}

\principle{\textit{\color{midgray}[TBD: Thesis to be written with Eric.]}}

\tbd[Principle body to be developed. Raw concept: [Eric to narrate the core idea].
This section is a placeholder---do not develop without Eric's direct input.]

\subsection{Story}
\tbd[Story to be attached.]

% ----------------------------------------------------------
\section{Principle 7}

\principle{\textit{\color{midgray}[TBD: Principle title and thesis to be confirmed with Eric.]}}

\tbd[Placeholder for Section I Principle 7.]

% ----------------------------------------------------------
\section{Principle 8}

\principle{\textit{\color{midgray}[TBD: Principle title and thesis to be confirmed with Eric.]}}

\tbd[Placeholder for Section I Principle 8.]
